\documentclass[12pt]{article}
\usepackage{array}
\usepackage{amsmath}
\usepackage{amssymb}
\usepackage[utf8]{inputenc}
\usepackage[russian]{babel}
\usepackage{geometry}
\usepackage{array}
\usepackage{multirow}
\usepackage{hyperref}

\geometry{a4paper, margin=2cm}

\title{Математическое обоснование симуляции времени маршрута}
\author{}
\date{}

\begin{document}

\maketitle

\section*{1. Описание модели}\mbox{}\\*

\subsection*{1.1 Сущности}\mbox{}\\*

Рассматриваемая схема метро состоит из ключевых сущностей:

\begin{itemize}
  \item \textbf{Станция} — $S$,
  \item \textbf{Переезд} — $R$,
  \item \textbf{Переход} — $T$.
\end{itemize}

Ветки не рассматриваются напрямую, их влияние косвенное.

\subsection*{1.2 Типы связей}\mbox{}\\*

Вычисления происходят между сущностями следующим образом:

\begin{table}[h]
\centering
\begin{tabular}{|c|c|c|c|}
\hline
Связь & Станция ($S_2$) & Переезд ($R_2$) & Переход ($T_2$) \\
\hline
Станция ($S_1$) & :x: & :white\_check\_mark: & :white\_check\_mark: \\
\hline
Переезд ($R_1$) & :white\_check\_mark: & :x: & :x: \\
\hline
Переход ($T_1$) & :white\_check\_mark: & :x: & :x: \\
\hline
\end{tabular}
\end{table}

Таким образом, необходимо рассчитать время для следующих типов переходов:

\begin{enumerate}
  \item $S \to R$,
  \item $R \to S$,
  \item $S \to T$,
  \item $T \to S$.
\end{enumerate}

\section*{2. Параметры}\mbox{}\\*

\subsection*{2.1 Локальные параметры}\mbox{}\\*

Локальные параметры задаются на этапе создания сущностей и остаются неизменными в рамках одной симуляции.

\subsubsection*{Параметры сущностей}\mbox{}\\*

\begin{enumerate}
\item \textbf{Станция (S)}:
  \begin{itemize}
    \item \texttt{Occupancy} — уровень загруженности станции;
  \end{itemize}
\item \textbf{Переезд (R)}:
  \begin{itemize}
    \item \texttt{Duration} — среднее время движения поезда по перегону;
  \end{itemize}
\item \textbf{Переход (T)}:
  \begin{itemize}  
    \item \texttt{Occupancy} — уровень загруженности перехода;
    \item \texttt{Duration} — среднее время движения пассажира по переходу при нулевой загруженности.
  \end{itemize}
\end{enumerate}

\subsubsection*{Ограничения локальных параметров}\mbox{}\\*

\begin{tabular}{|l|>{\raggedright\arraybackslash}m{7cm}|l|l|}
\hline
Имя & Описание & Тип & Ограничения \\
\hline
\texttt{Occupancy} & Уровень загруженности по x-бальной шкале & Целое число $\geq 0$ & $[0, 10]$ \\
\hline
\texttt{Duration} & Временной отрезок, секунды & Целое число $\geq 0$ & $[0, +\infty)$ \\
\hline
\end{tabular}

\subsection*{2.2 Глобальные параметры}\mbox{}\\*

\subsubsection*{Произвольные константы}\mbox{}\\*

\begin{itemize}  
  \item \texttt{AverStationEntryTime} или \texttt{asEntryT} или $\zeta$ — среднее время входа пассажира на станцию;
  \item \texttt{AverStationExitTime} или \texttt{asExitT} или $\eta$ — среднее время выхода пассажира со станции;
  \item \texttt{AverTrainWaitTime} или \texttt{atwaitT} или $\tau$ — среднее время ожидания поезда на станции;
  \item \texttt{QuintileTrustLevel} или \texttt{qtl} или $\gamma$ — уровень доверия квантилей.
\end{itemize} 

\subsubsection*{Характеристические данные}\mbox{}\\*

\begin{itemize}
  \item \texttt{AverChartBranchesStationsOccupancy} или \texttt{acbsOcc} - средняя загруженность станций веток схемы;
  \item \texttt{AverChartBranchesRailwaysDuration} или \texttt{acbrDur} - среднее время движения поезда в перегонах веток схемы;
  \item \texttt{AverBranchStationsOccupancy} или \texttt{absOcc} - средняя загруженность станций ветки;
  \item \texttt{AverBranchRailwaysDuration} или \texttt{abrDur} - среднее время движения поезда в перегонах ветки.
\end{itemize}

\subsubsection*{Ограничения глобальных параметров}\mbox{}\\*

\paragraph{Произвольные константы:}\mbox{}\\*

\begin{tabular}{|l|>{\raggedright\arraybackslash}m{5cm}|l|l|}
\hline
Имя & Описание & Тип & Ограничения \\
\hline
\texttt{AverStationEntryTime} & Среднее время входа на станцию, с. & Целое $\geq 0$ & $[0, +\infty)$ \\
\hline
\texttt{AverStationExitTime} & Среднее время выхода со станции, с. & Целое $\geq 0$ & $[0, +\infty)$ \\
\hline
\texttt{AverTrainWaitTime} & Среднее время ожидания поезда, с. & Целое $\geq 0$ & $[0, +\infty)$ \\
\hline
\texttt{QuintileTrustLevel} & Уровень доверия квантилей & Плавающее $\in [0, 1]$ & $[0, 1]$ \\
\hline
\end{tabular}

\paragraph{Характеристические данные:}\mbox{}\\*

\begin{tabular}{|l|>{\raggedright\arraybackslash}m{3.6cm}|l|l|}
\hline
Имя & Описание & Тип & Ограничения \\
\hline
\texttt{AverChartBranchesStationsOccupancy} & Средняя загруженность станций схемы & Целое $\geq 0$ & $[0, +\infty)$ \\
\hline
\texttt{AverChartBranchesRailwaysDuration} & Среднее время поездок в перегонах & Целое $\geq 0$ & $[0, +\infty)$ \\
\hline
\texttt{AverBranchStationsOccupancy} & Средняя загруженность станций ветки & Целое $\geq 0$ & $[0, +\infty)$ \\
\hline
\texttt{AverBranchRailwaysDuration} & Среднее время поездок по ветке & Целое $\geq 0$ & $[0, +\infty)$ \\
\hline
\end{tabular}

\section*{3. Вычисления}\mbox{}\\*

\subsection*{3.1 Локальные параметры}\mbox{}\\*

Пусть:

\begin{itemize}
  \item $S_i$ — $i$-я станция, где $i = 1, ..., S_N$, $S_N$ — количество станций,
  \item $R_j$ — $j$-й переезд, где $j = 1, ..., R_N$, $R_N$ — количество переездов,
  \item $T_k$ — $k$-й переход, где $k = 1, ..., T_N$, $T_N$ — количество переходов.
\end{itemize}

Тогда:

\begin{itemize}
  \item $\text{Socc}_i$ — загруженность $i$-й станции,
  \item $\text{Rdur}_j$ — среднее время поезда по $j$-му переезду,
  \item $\text{Tocc}_k$ — загруженность $k$-го перехода,
  \item $\text{Tdur}_k$ — среднее время движения по $k$-му переходу.
\end{itemize}

Все эти данные берутся из заготовленных таблиц.

\subsection*{3.2 Глобальные параметры}\mbox{}\\*

\subsubsection*{Произвольные константы}\mbox{}\\*

Задать вручную:

\begin{itemize}
  \item $\zeta$ --- \texttt{AverStationEntryTime},
  \item $\eta$ --- \texttt{AverStationExitTime},
  \item $\tau$ --- \texttt{AverTrainWaitTime},
  \item $\gamma$ --- \texttt{QuintileTrustLevel}.
\end{itemize}

\subsubsection*{Характеристические данные}\mbox{}\\*

\begin{itemize}
  \item \texttt{AverChartBranchesStationsOccupancy} или \texttt{acbsOcc} или $\mathcal{O}$:
    \begin{equation}
      \mathcal{O} = \frac{1}{R_N} \sum_{i=1}^{S_N} \text{Socc}_i
      \label{eq:characteristic:acbsOcc}
    \end{equation}
  \item \texttt{AverChartBranchesRailwaysDuration} или \texttt{acbrDur} или $\Upsilon$:
    \begin{equation}
      \Upsilon = \frac{1}{R_N} \sum_{j=1}^{R_N} \text{Rdur}_j
      \label{eq:characteristic:acbrDur}
    \end{equation}
  \item \texttt{AverBranchStationsOccupancy} или \texttt{absOcc} или $\mathrm{o}_b$:
    \begin{equation}
      \mathrm{o}_b = \frac{1}{N_b} \sum_{i}^{N_b} \text{Socc}_i
      \label{eq:characteristic:absOcc}
    \end{equation}
    для станций, принадлежащих $b$-ой ветке.
  \item \texttt{AverBranchRailwaysDuration} или \texttt{abrDur} или $\upsilon_b$:
    \begin{equation}
      \upsilon_b = \frac{1}{N_b} \sum_{j}^{N_b} \text{Rdur}_j
      \label{eq:characteristic:abrDur}
    \end{equation}
    для переездов, принадлежащих $b$-ой ветке.
\end{itemize}

\end{document}